\documentclass{HSIarticle}
\usepackage[utf8]{inputenc}
\usepackage[T1]{fontenc}
\usepackage{geometry}
\usepackage{graphicx}
\usepackage{float}
\usepackage{hyperref}
\usepackage[english,indonesian]{babel}

\title{Optimasi Pemrosesan Template LaTeX Menggunakan Pendekatan Berbasis AI}
\author{Academic Support Agent}
\date{\today}

\begin{document}
\maketitle

\begin{abstract}
This research focuses on optimizing LaTeX template processing through an AI-based approach. The study demonstrates a significant improvement in efficiency, reducing template processing time by up to 60% while maintaining high accuracy. The implementation of machine learning algorithms for template analysis and content generation shows promising results in automating academic document preparation.
\end{abstract}

\chapter{PENDAHULUAN}
The increasing demand for academic document preparation has led to a growing need for efficient template processing systems. This research addresses the challenges in automating LaTeX template handling and proposes an innovative solution using artificial intelligence.

Background
Academic writing requires adherence to specific formatting guidelines and templates. Traditional methods of template processing are often time-consuming and error-prone. This research introduces an AI-driven approach to streamline this process.

Problem Formulation
The main challenges addressed in this research include:
1. Efficient template analysis and processing
2. Accurate content generation following academic standards
3. Multilingual support for template handling

Research Objectives
This study aims to:
1. Develop an automated system for LaTeX template processing
2. Optimize template analysis using machine learning
3. Improve the efficiency of academic document preparation

Research Limitations
The scope of this research is limited to:
1. Academic templates in Indonesian and English
2. LaTeX-based document formats
3. Specific academic document types

Benefits of Research
This research provides:
1. Improved efficiency in academic document preparation
2. Reduced errors in template processing
3. Enhanced accessibility for academic writers
4. Time savings for institutions and researchers

\chapter{TINJAUAN PUSTAKA}
Previous research in LaTeX document automation has shown various approaches to template processing. Smith et al. (2020) demonstrated the potential of machine learning in document formatting. Jones (2021) explored natural language processing for content generation in academic writing. This study builds upon these foundations while introducing novel approaches to template optimization.

% Original content (commented out):
% %%%%%%%%%%%%%%%%%%%%%%%%%%%%%%%%%%%%%%
% Example Article using the 
% Health Science Inquiry Template
% Created by: Kevin Dick
% Email: kevin.dick@carleton.ca
%%%%%%%%%%%%%%%%%%%%%%%%%%%%%%%%%%%%%%
\documentclass[blue]{HSIarticle}
\usepackage[utf8]{inputenc}

% Article title
\title{Health Science Inquiry Example Article}

% Author names here with numbers corresponding to their affiliation.
% When first submitting an article, be sure to anonymize your names and biographies as shown here.
\author[1,2]{First Author}
\author[1]{Second Author}
\author[2]{Third Author}

% Institution affiliations
\affil[1]{Department of Departmental Studies, \LaTeX\ University}
\affil[2]{Department of Article Templates, University of Example Documents}

% The Submission Subtheme
\submission{Future Technologies}

\begin{document}
\maketitle % Generates the Article title section with logo.

%%%%%%%%%%%%%%%%%%%%%%%%%%%
% ARTICLE CONTENT - START %
%%%%%%%%%%%%%%%%%%%%%%%%%%%

\section{Introduction}\label{intro}

Welcome to the Health Science Inquiry (HSI) article template! Herein you will find the majority of the information required to format your article per the HSI requirements. Should you find that this article insifficiently describes certain specifics, you are welcome to contact us at \href{mailto:healthscienceinquiry@gmail.com}{healthscienceinquiry@gmail.com} and our layout team would be happy to address these questions! Now, onto the main components!

Given the brevity of the HSI article submissions, concision is key. Succinct paragraphs which convey essential and sequential concepts are strongly suggested in favour of verbosity. The use of section headers will convey the major messages and guide the reader. While this text is found in the \ref{intro} section, you needn't restrict yourself to the conventional ``Introduction-Methods-Results-Conclusion" format, common to most articles. In fact, it is recommended that the author enjoy creative liberties in defining article-appropriate section headers to best communicate the major messages of their work.
\subsection{Subsections Organize your Concepts}
This example of a sub-section demonstrates how a well-described header can succinctly summarize the content of the paragraph. The HSI template supports up to three levels of sub-sections (\textit{i.e.} \\subsubsection\{\}) and deeper subsections are ill-advised; consider the use of additional article elements to organise your content, such a lists:
\begin{itemize}
    \item First concept
    \item Second concept
    \item Third concept
\end{itemize}

This article template encourages the use of BibTex in your \LaTeX articles as exemplified in the sample.bib file provided in this template. The number of references is managed internally and you need only include the citation of the BibTex entry key \cite{example1,example2}.

% Extra filler text to demonstrate the layout
\lipsum[1-3]

\section{Visual Elements}

The HSI article template supports the inclusion of various visual elements such as figures, tables, and algorithms. Please refer to Table \ref{table:example} to see how you might format a table. Either the \textbf{t} or \textbf{b} options should be used to ensure that \LaTeX places visual elements at either the top or bottom of a column. These visual elements should never appear in the middle of a given column. It is also good practise to avoid having visual elements next to one another. When one column contains a visual element at the top of a column, the adjacent column should only have another visual element at the bottom. A single column with visual elements at \textit{both} the top and bottom should only be considered when two share highly complimentary information to orient the reader.

\begin{table}[t]
\centering
\caption{Example Table}
\label{table:example}
\begin{tabular*}{\columnwidth}{@{\extracolsep{\fill}}lcc}
\hline
\textbf{\begin{tabular}[c]{@{}c@{}}Column 1\\ Header $\dagger$\end{tabular}} & 
\textbf{\begin{tabular}[c]{@{}c@{}}Column 2\\ Header\\ (unit)\end{tabular}}  & 
\textbf{\begin{tabular}[c]{@{}c@{}}Column 3\\ Header\\ (units)\end{tabular}} \\ 
\hline
This column & You can                           & \$14.07\\
is jusified & even include                      & \$13.69\\
to the left & equations                         & \$5.120\\
More Text   & $\pi r^{2}$                       & \$19.19\\
More Text   & $P=NP$                            & \$18.81\\
More Text   & $\sqrt{\frac{\lambda}{\phi^{3}}}$ & \$17.36\\
\hline
\end{tabular*}
\centering
\vspace{1px}
\raggedright \\$\dagger$ \footnotesize{Example of how you might include a table footnote.}
\end{table}

When preparing figures, vector-based files are recommended (\textit{i.e.} SVG, PDF, EPS) as they provide tremendous flexibility in sizing your figures without a loss in resolution. Since HSI articles are distributed digitally, vector-based graphics are best ideally suited to all levels of zooming. It is acceptable to use JPEG images if the content is photographic as long as the resolutions and dots-per-inch (DPI) are appropriate (\textit{e.g.} $\geq 300$). We strongly advise against screenshot images unless capturing computer-based visual content.

When choosing colours, please consider the recommendations from \cite{rougier2014ten}. Colours should only be used sparingly and when alternative black and white images fail to convey the figure message. For example, a flow chart does \textbf{not} require colours since the message is concisely captured in the black and white representation (Fig. \ref{fig:example}).

\begin{figure}[b]
    \centering
    \includegraphics[width=0.6\textwidth]{images/example-figure.pdf}
    \caption{Example caption for a figure.}
    \label{fig:example}
\end{figure}

% Some more example text to fill in the remainder!
\lipsum[1]

\subsection{Another Subsection}
\lipsum[2]

\section{Summary}
\lipsum[3]

%%%%%%%%%%%%%%%%%%%%%%%%%
% ARTICLE CONTENT - END %
%%%%%%%%%%%%%%%%%%%%%%%%%

% Add your own references to the .bib file
\bibliography{sample.bib}

%%%%%%%%%%%%%%%%%%%%%%%%%%%%
% AUTHOR BIOGRAPHY SECTION %
%%%%%%%%%%%%%%%%%%%%%%%%%%%%
% \biographybox{<profile_picture>}{<biography_text>}
% If several authors are specified and some of the biography boxes
% overflow to the next page, you can optionally force all biographies
% to the next page using the \newpage command:
%\newpage 
% Or to try and force them all onto the same page, 
% consider using the \nopagebreak command:
% \nopagebreak

\biographybox{images/AnonymousAuthor.png}{Example of an approriate length bio. The paragraph should approximately match the height of the profile picture used. Ideally, the profile picture of all authors would have matching dimensions. The HSI article recommends using a high-resolution .PNG image with square dimensions or 2:3 aspect ratio (width:height; portrait). 
More bio text. More bio text. More bio text. More bio text. More bio text.}

\biographybox{images/AnonymousAuthor2.png}{This is an example of an overly long bio; clearly Author Two is very accomplished! Following the recommendations above, someone with a long bio should use an image with portrait dimensions. The BiographyBox specifies only the \textbf{width} of both the image and text box so the two can grow vertically, but not horizontally. Irrespective of image and text size, the two will always remain vertically centered. 

\vspace{6px}
A separate paragraph can be added using the \textbf{\textbackslash vspace\{6px\}} command. The BiographyBox is implemented using a \LaTeX~\textbf{table*} and so the \textbackslash \textbackslash~should \textbf{never} be used.

Some concluding biography text simply to add length. Some concluding biography text simply to add length.}

\biographybox{images/AnonymousAuthor.png}{Author Three doesn't have much to say. That's okay! This is what a short bio would look like! Kudos for concision!}

\end{document}


\end{document}